\section{Introduction}
\label{sec:introduction}

\subsection{Brief Introduction to Options and Implied Volatility}
\label{subsec:ivs}

Option contracts are one of the most important types of derivative contracts in the context of financial markets.
A derivative is a contract traded with respect to another asset (also called the underlying asset)
and an option contract allows the buyer of the contract to sell or buy the specified asset from the seller at a predefined future date.
The pricing equations for these derivatives, called the Black-Scholes model, require knowing the values for the time to maturity of the contract,
the price of the underlying asset to be traded, the price of the asset to be traded and the market's risk-free rate at the time~\citep{Black1973}.
An additional variable, called the implied volatility (\textit{IV} or $\sigma_{iv}$) is also required for the pricing of these contracts.
This value refers to the market's expected value for the volatility of the asset price up until the time of maturity.
% and the implied volatility surfaces play a critical role in finding the correct no-arbitrage price~\cite{Ludkovski2023} of these contracts.

Having access to the implied volatility values allows investors to make sales at the prices consistent with the current expected contract price,
that is, the no-arbitrage price according to the market expectancy.
However, the \textit{IV} values corresponding to all possible strike prices and maturities are not directly observable in the market.
Ensuring a way to obtain these values with confidence and readily available data in the market is a crucial task for quantitative analysts~\citep{Kim2006}.
Usual strategies for estimating these values daily are predominantly based on linear or spline interpolation
algorithms~\citep{Homescu2011} or stochastic differential equation models~\citep{Jordan2011}
that use numerical methods to approximate the values from existing listed contracts.
Both methods try to approximate the correct value for the implied volatility for each pair of time to maturity and strike price.
This approach allows visualizing the model results in the form of a surface,
where the \textbf{XY} axes are the maturity and strikes pairs and the \textbf{Z} axis is the \textit{IV}.
This surface, called the \textbf{implied volatility surface} (IVS), noted as $\mathcal{S}$ is therefore a function of the maturity time $t$ and the strike price $K$,

\begin{equation}
    \label{eq:ivs}
    \mathcal{S} \coloneqq f(t, K), \quad \text{where } f \colon \mathbb{R}^{2} \to \mathbb{R}.
\end{equation}

Ensuring accurate predictions of the IVS is crucial, as even small errors can lead to incorrect pricing and consequently
large losses of profit or even negative return~\citep{Hagan2014} given the leveraging capabilities of options when compared to non-derivative securities.
Classic models such as linear interpolation, polynomial fitting, and \textit{splines} not only lack flexibility for extrapolation,
but also generate either non-smooth surfaces, arbitrage opportunities, or both~\citep{Hagan2014}.

\subsection{Contributions of the Presented Work}
\label{subsec:contributions}

In this paper, we propose an approach to address the challenges of interpolation and extrapolation of the IVS, focusing on reducing arbitrage opportunities.
The chosen method uses a stochastic differential equation solution approach, specifically the \textit{Stochastic Alpha, Beta, Rho} (\textit{SABR}) model.
The parameters $\alpha, \rho, \nu$ of the \textit{SABR} model are obtained through the
Broyden-Fletcher-Goldfarb-Shanno (BFGS) quasi-Newton non-linear optimization algorithm~\cite{Broyden1970, Fletcher1970, Goldfarb1970, Shanno1970}.
This step is commonly performed using only a subset of the points observed daily~\cite{Kim2021},
due to the fact that using all points introduces various outliers that do not match the real prices of the derivative, which in turn increases the variance of the final model.
To solve this difficulty, the points chosen for the optimization step are often manually filtered according to criteria set by human experts~\cite{Hagan2002},
and only a candidate subset of the original points is used.
Rather than manually selecting these candidate points, the use of models capable of receiving point clouds and returning a score, such as transformers, to select the best points
automatically has been proposed by~\citet{Hagan2002} instead.
The main contribution of this research will be the adaptation of \textit{neural implicit layers},
a physics-inspired variation of fully connected linear layers that outputs a solution to a fixed point equation, to be used as the output layer to a transformer architecture.
This approach is an implicit fixed-point solver variation for the scoring model proposed by~\citet{Kim2021}, aiming to enforce a zero arbitrage condition in the optimization process.
This type of layer has efficiently been shown to perform well representing smooth surface densities and demonstrated state-of-the-art performance in function
fitting tasks similar to the one presented here~\cite{Kolter2020, Kawaguchi2021},
all while being computationally faster for solving interpolation tasks than linear layer models with similar performance~\cite{Kolter2020}.
Additionally, transformers have also been shown to perform better at handling unordered point clouds over time compared to recurrent neural layers~\citep{Kim2024}.

\subsection{The Transformer Architecture}
\label{subsec:transformer}
The aforementioned \textit{SABR} model is a stochastic volatility model used to obtain smooth surfaces, i.e.,
continuous and infinitely differentiable at all points~\cite{Hagan2015}.
In the context of this work, a score will be assigned for each candidate point of the daily listed options contracts,
called the \textit{summarized suitability} score (\textit{SS} score) by \citet{Kim2021}.
This score corresponds to the suitability of the candidate point to summarize the characteristics of the generated surface, i.e.,
how feasible it is to reduce the number of candidates when optimizing the values for one of the parameters of the SABR model
without altering the resulting surface shape or its smoothness\footnote{The use of a per-point score is similar to what the principal components of a collection of points is when using the Principal Component Analysis (PCA) technique.}.
In summary, the score is used to select the highest valued points, and therefore the most representative of the surface,
removing those that might introduce arbitrage in the optimization process for each of the \textit{SABR} models~\cite{Kim2021}
(a separate model is fitted for each unique strike price in a given day).
\hyperref[subsec:sabr]{Section~\ref*{subsec:sabr}}
covers the details of fitting the \textit{SABR} model at each strike price.

\subsection{Fixed-Point problems}
\label{subsec:implicit}

Differently from the usual function approximation tasks, a fixed point problem is instead defined by conditions that the output must satisfy.
For explicit functions, a given output $z$ is computed directly from the input $x$ using a function $f$, such that $z = f(x)$.
In contrast, an implicit function  $g$ is such that it depends on both the input $x$ and the expected output $z$,
where $z$ is a solution for the equation:

\begin{equation}
    \label{eq:fixed_point}
    g(x, z) = 0
\end{equation}

Formulating the problem as a fixed point equation can be used for a multitude of problems, including finding fixed points in recurrent networks,
solutions to differential equations leading to Neural ordinary differential equations, and several physics optimization problems~\cite{Chen2019}.
The implicit formulation offers several practical advantages, primarily the separation of the solution method from the layer
definition which can be useful in maintaining interpretability of the resulting model\cite{Kawaguchi2021} and as a way to enforce constraints in the optimization process.

Moreover, the resulting implicit layer offers significant benefits specifically in deep learning and automatic differentiation frameworks~\cite{Massaroli2021,Kolter2020}.
Automatic differentiation methods store the entire computation graph, which can be memory-intensive.
With implicit layers, however, the implicit function theorem allows computing gradients directly at the solution point and
therefore storing intermediate variables becomes redundant~\cite{Massaroli2021}.
This improves memory efficiency, making implicit layers particularly advantageous when substituting existing layers within large deep learning models~\cite{Li2020}.

\subsection{Anderson Acceleration for Implicit Layers}
\label{subsec:anderson}

The Anderson Acceleration (AA) method is a technique used to accelerate the convergence of fixed-point iterations.
It is particularly useful when the fixed-point iteration is slow to converge or when the fixed-point iteration is used in a deep learning context.
The method is based on the idea of extrapolating the current fixed-point iteration using a linear combination of the previous iterations.
For this work, we use a Equilibrium Model variation of the Anderson Acceleration method, which is particularly useful for implicit layers in neural networks.
The AA method is used to solve the fixed-point equation within the layer, by extrapolating the current fixed-point iteration using a linear combination of the previous iterations.
The AA method is particularly useful for implicit layers in neural networks, as it allows for faster convergence of the fixed-point iteration.
Our proposed architecture uses the AA method to solve the fixed-point equation in the implicit layer,
according to the \hyperref[alg:anderson]{Algorithm~\ref*{alg:anderson}} for our implementation of the AA method.

\begin{algorithm}
    \begin{algorithmic}[1]
        \Require Function \( f \), initial point \( x_0 \), memory size \( m \), damping factor \( \lambda \), maximum iterations \( K \), tolerance \( \epsilon \), mixing factor \( \beta \)
        \State Initialize \( X_0 = x_0, X_1 = F_0 = f(x_0), F_1 = f(F_0) \)
        \For{$k = 2$ to \( K \)}
            \State \( j \gets \min(k, m) \)
            \State \( G = F_{0, \ldots, j} - X_{ 0, \ldots, j} \)
            \State \( H = G G^\top + \lambda I \)
            \State Solve \( H^{-1} \cdot y = \alpha \)
            \State \( \alpha \gets \alpha_{1, \ldots, n + 1} \)
            \State \( X_k \gets \beta \alpha^\top F_0 + (1 - \beta) \alpha^\top X_0 \)
            \State \( F_k \gets f(X_k) \)
            \State Compute residual:
            \[
                \text{residual} = \frac{\| F_k - X_k \|}{\| F_k \| + 10^{-5}}
            \]
            \If{residual \(<  \epsilon \)}
                \State \textbf{break}
            \EndIf
        \EndFor
        \State \Return \( X_k \)
    \end{algorithmic}
    \caption{Equilibrium Model with Anderson Acceleration}
    \label{alg:anderson}
\end{algorithm}

The final deep equilibrium model used within the implicit layer uses the Anderson Acceleration algorithm to solve the fixed-point equation,
and the backpropagation algorithm is used to find the weights that minimize the error function of the model.
The solution to the fixed-point equation, is then used as the output of the layer, as shown by \hyperref[alg:forward]{Algorithm~\ref*{alg:forward}}.

\begin{algorithm}
    \begin{algorithmic}[1]
        \Function{ImplicitLayer}{$x, m, \lambda, \epsilon, \beta$}
            \State \textbf{function} $f(x) \gets W \cdot x + b$
            \State $x \gets \texttt{AA}(f, x, m, \lambda, \epsilon, \beta)$
            \State $y \gets x \cdot W^{T} + b$

            \State \Return $y$
        \EndFunction
    \end{algorithmic}
    \caption{Implicit layer forward pass with Anderson acceleration (AA).}
    \label{alg:forward}
\end{algorithm}
